\documentclass[11pt]{article}

\usepackage[T1]{fontenc}
\usepackage{lmodern}
\usepackage{amsmath,amssymb,amsthm}
\usepackage[margin=1in]{geometry}
\usepackage{microtype}
\usepackage[hidelinks]{hyperref}
\hypersetup{
  pdftitle={A 2-adic obstruction for an eighteen-variable rational moment system},
  pdfauthor={Xiaojian Zeng}
}

\newtheorem{theorem}{Theorem}
\newcommand{\Q}{\mathbb{Q}}

\title{A 2-adic obstruction for an eighteen-variable\\rational moment system}
\author{Xiaojian Zeng\\
  \small PBC School of Finance, Tsinghua University, 100190\\
  \small \texttt{xiaojianzeng@foxmail.com}}
\date{}

\begin{document}
\maketitle

\begin{abstract}
We prove that there is no rational \(18\)-tuple whose second, fourth, and
sixth power sums are respectively \(9\), \(27/4\), and \(45/8\).  No
distinctness, nonzero-coordinate, sign, or interval hypothesis is imposed.
After clearing a primitive common denominator, congruences modulo \(32\)
exclude every possible 2-adic valuation of that denominator.  The proof is
elementary and does not depend on a computational search.
\end{abstract}

\section{Context and statement}

For a probability measure, a Chebyshev-type quadrature is an equal-weight
formula matching a prescribed initial segment of moments~\cite{Peled2010}.
In the interval-design literature, the qualifier ``strict'' requires
distinct nodes, whereas
the unqualified notion allows repetitions~\cite{MisawaMunemasaSawa2026}.
The even moments of the arcsine probability measure on \([-1,1]\) are
\[
  \int_{-1}^{1}t^{2r}\frac{dt}{\pi\sqrt{1-t^2}}
  =\frac{\binom{2r}{r}}{4^r}.
\]
For the corresponding Jacobi weight
\(w_{-1/2,-1/2}\), the Gaussian rule has equal weights, and an
\(N\)-node Gaussian rule is exact through degree \(2N-1\)
~\cite{GilboaPeled2017}.  Taking \(N=18\) therefore gives a real
\(18\)-tuple with the three target moments below.  The obstruction proved in
this paper is specifically to choosing all eighteen coordinates rational.
Consequently the normalized target moments \(1/2,3/8,5/16\) lead, for
eighteen equal-weight rational coordinates, to the arithmetic system below.
These are precisely the Hilbert--Kamke equations attached to antipodal
rational \(7\)-designs for the Chebyshev measure; Mishima--Lu--Sawa--Uchida
exclude \(N\leq16\) and pose the range \(N\geq17\)~
\cite{MishimaLuSawaUchida2026}.  The system by itself is weaker than a strict
quadrature or design problem: it does not require nodes in \([-1,1]\),
distinct nodes, or the matching of odd moments.

A closely related preprint concerns the analogous \(17\)-variable system;
its deposited abstract reports nonexistence, also without distinctness or
interval assumptions, by an elementary 2-adic argument~
\cite{Zeng2026SSRN7385178}.  The
proof below independently treats all eighteen variables and uses only the
displayed moment equations.  We make no priority claim beyond this comparison.

\begin{theorem}\label{thm:main}
There is no tuple \((x_1,\ldots,x_{18})\in\Q^{18}\) satisfying
\[
 \sum_{j=1}^{18}x_j^2=9,\qquad
 \sum_{j=1}^{18}x_j^4=\frac{27}{4},\qquad
 \sum_{j=1}^{18}x_j^6=\frac{45}{8}.
\]
This assertion allows zero and repeated coordinates, arbitrary signs, and
coordinates outside \([-1,1]\).
\end{theorem}

\section{The congruence obstruction}

\begin{proof}[Proof of Theorem~\ref{thm:main}]
Suppose that a solution exists.  Choose integers
\(a_1,\ldots,a_{18}\) and a positive integer \(D\) such that
\[
  x_j=\frac{a_j}{D},\qquad
  \gcd(D,a_1,\ldots,a_{18})=1.
\]
Writing \(S_r=\sum_{j=1}^{18}a_j^r\), clearing denominators gives
\begin{equation}\label{eq:cleared}
 S_2=9D^2,\qquad S_4=\frac{27D^4}{4},\qquad
 S_6=\frac{45D^6}{8}.
\end{equation}
The denominator \(D\) is even.  Indeed, if \(D\) were odd, the identity
\(4S_4=27D^4\) would equate an integer divisible by \(4\) with an odd
integer.  Put \(D=2^k d\), where \(k\geq1\) and \(d\) is odd.  Primitivity
then ensures that at least one \(a_j\) is odd.

Let \(O\) denote the number of odd \(a_j\).  For every odd integer \(a\),
define \(p(a)\in\mathbb Z/4\mathbb Z\) by
\[
 a^2\equiv1+8p(a)\pmod {32}.
\]
Squaring and multiplying this congruence gives
\begin{equation}\label{eq:odd-powers}
 a^4\equiv1+16p(a),\qquad
 a^6\equiv1-8p(a)\pmod {32}.
\end{equation}
Set
\[
 P=\sum_{a_j\ \mathrm{odd}}p(a_j)\in\mathbb Z/4\mathbb Z,
\]
and let \(E\) and \(F\) count the indices for which \(v_2(a_j)=1\) and
\(v_2(a_j)=2\), respectively.  Modulo \(32\), an integer of 2-adic
valuation \(1\) contributes \((4,16,0)\) to
\((a^2,a^4,a^6)\); one of valuation \(2\) contributes \((16,0,0)\); and
one of valuation at least \(3\) contributes \((0,0,0)\).  Together with
\eqref{eq:odd-powers}, this yields
\begin{align}
 S_2&\equiv O+8P+4E+16F,\label{eq:s2}\\
 S_4&\equiv O+16P+16E,\label{eq:s4}\\
 S_6&\equiv O-8P\pmod {32}.\label{eq:s6}
\end{align}

First suppose \(k=1\).  Because \(d\) is odd, the right-hand sides of
\eqref{eq:cleared} are congruent to \((4,12,8)\) modulo \(32\).  Reducing
the fourth-power identity modulo \(16\) gives
\(O\equiv12\pmod {16}\).  Since \(1\leq O\leq18\), this forces \(O=12\).
The sixth-power identity modulo \(8\), however, gives
\(O\equiv0\pmod8\), a contradiction.

It remains to consider \(k\geq2\).  In this case both \(S_4\) and \(S_6\)
vanish modulo \(32\).  Equation~\eqref{eq:s4} modulo \(16\) gives
\(O\equiv0\pmod {16}\), and hence \(O=16\).  Equation~\eqref{eq:s6}
then gives \(P\equiv2\pmod4\), while substitution into
\eqref{eq:s4} shows that \(E\) is odd.  Exactly two numerators are even,
so \(E=1\).  The other even numerator has valuation \(2\) or at least
\(3\), and therefore \(F\in\{0,1\}\).  Equation~\eqref{eq:s2} now gives
\[
 S_2\equiv4+16F\in\{4,20\}\pmod {32}.
\]
On the other hand, the first identity in \eqref{eq:cleared} gives
\(S_2\equiv16\pmod {32}\) when \(k=2\), and
\(S_2\equiv0\pmod {32}\) when \(k\geq3\).  Both possibilities are
contradictions.  Thus every possible value of \(v_2(D)\) is excluded.
\end{proof}

\section{Exact verification and limitations}

The residue calculation admits a finite exact check that is independent of
the case analysis.  For each residue \(a\in\mathbb Z/32\mathbb Z\), record
\((a^2,a^4,a^6,\epsilon(a))\), where \(\epsilon(a)\) records whether \(a\)
is odd.  Starting from the zero state, iteratively add one such contribution
for each of eighteen coordinates, always reducing the first three entries
modulo \(32\).  Exact integer dynamic programming leaves \(520\) compressed
states.  Among states with an odd coordinate, none has power-sum triple
\[
 (4,12,8),\qquad (16,0,0),\qquad (0,0,0),
\]
the respective targets for \(v_2(D)=1\), \(v_2(D)=2\), and
\(v_2(D)\geq3\).  This finite check corroborates the proof but is not used
to establish the theorem.

The result resolves only the rational system in Theorem~\ref{thm:main}.
It does not assert nonexistence over the reals, nor does it classify systems
with other numbers of coordinates or other target moments.

\paragraph{Computational and AI assistance.}
AI tools assisted with proof search, exact finite-check design, audit support,
and manuscript preparation.  The modular proof was subsequently reconstructed
and checked line by line in a separate mathematical review.  The dynamic
program described above used exact integer arithmetic and supplies corroboration
only; no floating-point calculation or computational search is an input to the
proof.  Bibliographic claims about the \(17\)-variable preprint are restricted
to its deposited metadata and abstract.

\bibliographystyle{abbrv}
\bibliography{references}

\end{document}
